\section{Conceptos}

Los cuaterniones números hipercomplejos, son una extensión de los números complejos hacia un espacio de cuatro dimensiones, compuestos por una parte real y tres partes imaginarias distintas ($i$, $j$ y $k$).

Un cuaternión se puede ver así:

\begin{equation}
    a + bi + cj + dk
\end{equation}

donde $i$, $j$ y $k$ son tres números imaginarios distintos; $a,b,c,d\in\mathbb{R}$; y los términos con $i$, $j$ y $k$ son las tres partes imaginarias distintas.

Si te preguntas, cuánto vale exactamente $i$, cuánto vale exactamente $j$ y cuánto vale exactamente $k$ en un cuaternión; la respuesta es que $i$, $j$ y $k$ ya son los valores exactos. Son elementos base irreductibles de este sistema algebraico de cuaterniones. En cuaterniones, $i\neq\sqrt{-1}$, sólo $i=i$, no se puede reducir a más, tampoco $j$ y $k$ se pueden reducir más. 

Existe una relación muy importante entre esos 3 números imaginarios que define a los cuaterniones, y esa relación es la siguiente:

\begin{equation}
    i^{2}=j^{2}=k^{2}=ijk=-1
\end{equation}

Ojo. La ecuación anterior no implica que:

\begin{equation}
    i=j=k=\sqrt{-1}
\end{equation}

Pues, como se mencionó anteriormente, $i$, $j$ y $k$ no se pueden reducir más, no son variables que toman un valor en el conjunto de los números reales ni imaginarios. $i$, $j$ y $k$ son los llamados números \textbf{hipercomplejos}, un cuaternión es un número hipercomplejo. 

Ahora bien, a pesar de que $i$, $j$ y $k$ no se pueden reducir a más, lo que sí podemos hacer es representarlos usando vectores de cuatro dimensiones. Podemos representarlas entonces como sigue:

\begin{align}
    i &= (0,1,0,0) \\
    j &= (0,0,1,0) \\
    k &= (0,0,0,1)
\end{align}

Y si quisiéramos representar a un número real (digamos, el número 1) como un vector de cuatro dimensiones también, entonces lo representaríamos así:

\begin{equation}
    1 = (1,0,0,0)
\end{equation}

Usamos los cuaterniones porque, en robótica y en otras áreas (videojuegos y modelado 3D), nos permiten hacer \textbf{rotaciones} de objetos dentro de un espacio tridimensional. Al ser los cuaterniones objetos matemáticos de 4 dimensiones, nos permiten hacer estas rotaciones de forma mucho \textbf{más eficiente computacionalmente}. Al aplicar un cuaternión en sistemas de coordenadas, las partes imaginarias $i$, $j$ y $k$ las interpretamos como coordenadas en los ejes $x$, $y$ y $z$ en un espacio tridimensional. Es decir, aunque los cuaterniones son objetos de 4 dimensiones, les damos una aplicación en nuestro mundo de 3 dimensiones. 

\subsection{Representaciones de un Cuaternión}

Entonces, un cuaternión es un llamado número \textbf{hipercomplejo}, pues consiste de una parte real y tres partes imaginarias distintas. Nótese que un cuaternión tiene 4 componentes. 

Luego entonces, podemos representar a un cuaternión de las siguientes formas:

\begin{itemize}
    \item Como un \textbf{Vector de Cuatro Componentes}. 
    \begin{align*} 
        (a,b,c,d) 
    \end{align*}

    \item Como una \textbf{Extensión de los Números Complejos}, como un \textbf{Número Hipercomplejo}, es decir, muy similar a como se representan los números complejos.
    \begin{align*} 
        a+bi+cj+dk
    \end{align*}
    
    \item Como la \textbf{Unión de un Escalar y un Vector con Tres Componentes}. 
    \begin{align*} 
        &a+(b,c,d)\\
        &=a+\vec{z}
    \end{align*}
    con $\vec{z}=(b,c,d)$
    
    \item Como un \textbf{Par ordenado de partes real e imaginaria}. 
    \begin{align*} 
        &(a,(b,c,d))\\
        &=(a,\vec{z})
    \end{align*}
    con $\vec{z}=(b,c,d)$
\end{itemize}

Nótese que todas las notaciones anteriores son representaciones diferentes \textbf{de un mismo} objeto matemático con 4 componentes.

Nótese también que en todas las representaciones anteriores que tienen el símbolo $+$, el símbolo $+$ no es una suma aritmética (e.g. no puedes sumar un escalar con un vector), es una \textit{suma formal}. Una \textit{suma formal} es una expresión que utiliza el símbolo de suma ($+$) para agrupar elementos (como símbolos, vectores, monomios o series) como una estructura algebraica, sin que exista necesariamente una operación numérica de adición calculable entre ellos. Es exactamente el mismo principio que se usa en los números complejos ($z = a + bi$). El operador $+$ no pide que sumes $a$ con $bi$, sino que los mantiene juntos como las dos dimensiones independientes de un solo número complejo. En el cuaternión, el $+$ separa la dimensión real de las tres dimensiones imaginarias/vectoriales.

Nota: En cuaterniones, el número $i$ es un cuaternión cuya parte real es cero y cuyas partes imaginarias $j$ y $k$ también son cero, excepto la parte con $i$, que vale 1.
\begin{equation*}
    i=0+1i+0j+0k
\end{equation*}

\subsubsection{Notación}
Conociendo las diferentes formas en las que se puede representar un cuaternion mostrando sus 4 componentes, también podemos representar a un cuaternión con un sólo símbolo como sigue:
\begin{equation}
    \tilde{q}
\end{equation}
Es decir, podemos representar a un cuaternion (sus cuatro componentes) con una letra con virgulilla.

Entonces tenemos:
\begin{equation}
    \tilde{q}=(q_0, q_1, q_2, q_3)=q_0+q_1i+q_2j+q_3k=q_0+\vec{q}=(q_0,\vec{q})
\end{equation}
donde $\vec{q}=(q_1,q_2,q_3)$.




